\begin{center}
    \large
    \begin{singlespace}    
        \textbf{\chineseTitle{}} \\[0.5cm]
    \end{singlespace}

    \begin{singlespace}    
    學生：\studentCnName  \hspace{2.5cm}  指導教授：\advisorCnName \hspace{0.1cm} 博士 \\

    \ifdefined\advisorCnNameB % 如果有共同指導教授
        \hspace{9.6cm}  \advisorCnNameB ~ 博士 \\
    \fi
    \end{singlespace}

    \vspace{0.5cm}

    \begin{singlespace}
        國立陽明交通大學\NameofDepartmentInstituteCN\\[0.2cm]
    \end{singlespace}

    \textbf{摘\hspace{1cm}要} \\[0.5cm]

\end{center}

\normalsize

知識圖譜感知（KG-aware）推薦系統在 KG 稠密且品質良好時普遍能提升準確度，但當 KG 稀疏或品質不佳時（即評論衍生、冷啟動、隱私受限等常見場景），既有研究探討有限。本文以書評衍生 KG 為例（過濾後每件商品平均 2.4 條邊），觀察到四個代表性方法（KGAT、KGCL、MCCLK、KGRec）皆不及完全不使用 KG 的純 LightGCN，原因在於這些方法將 KG 訊號內嵌於計分管線，無結構性開關可在 KG 不可信時將其關閉，KG 雜訊因此污染協同表示之梯度流。

本文提出 RA-GARK（Rationale-Aware Gating over Review-Aspect Knowledge graphs），將 KG 視為可閘控之側通道而非計分管線的必要組件，並由三項設計選擇實現。其一，本地偏置融合閘起始關閉，使訓練起點等同於純 LightGCN，僅在梯度顯示 KG 有助於協同訊號時才逐步開啟，提供結構性優雅退化。其二，softmax aspect-saliency 注意力將每件商品的 KG 語意彙整為少數潛在面向槽，並為給定的使用者—物品配對選取最具代表性者。其三，以 SVD 為基礎的初始化使面向槽自 KG 的語意幾何出發，而非起於隨機雜訊。在 KG 不變的條件下，RA-GARK 於 NDCG@20 達 0.1238，較最強 KG-aware 基準 KGRec 提升 13.1\%、較純 LightGCN 提升 5.0\%，在不更動 KG 的前提下將其貢獻自「無益」翻轉為「有益」。上述結果取自單一稀疏基準，顯示當 KG 稀疏或不可靠時，將 KG 視為可閘控之側通道是一個有效的設計原則。\\[0.5\baselineskip]
\textbf{關鍵字：}知識圖譜推薦；可退化機制；融合閘；稀疏知識圖譜；協同過濾
